\chapter{Zugänge \& Systemarchitektur}

\section{Zugangsübersicht}

\Sidebar{Sicherheit}
Die folgende Übersicht zeigt die wichtigsten Zugänge und Dienste, die für den Aufbau und Betrieb der Smart-Home-Lernumgebung relevant sind. Die Logins sind hier nur als Beispiel dargestellt und müssen vor der Nutzung durch die tatsächlichen Zugangsdaten ersetzt werden.

\begin{tabularx}{\linewidth}{X X X X >{\raggedright\arraybackslash}p{1cm} >{\raggedright\arraybackslash}p{3.6cm}}
    \rowcolor{SmartLight}
    \textbf{Bereich} & \textbf{Zugriff IP/ID} & \textbf{Port} & \textbf{Benutzer} & \textbf{Passwort} & \textbf{Hinweis} \\
    \hline
    WLAN  & MMSE-Smart-Home-Wifi & -- & -- & *** & Verbindung für Router, Server und Lernmodule \\
    \hline
    Router Config & 10.0.0.138 & 80 / 443 & admin & *** & lokale Administration im Projekt-Netz \\
    \hline
    SSH & 10.0.0.1/\textit{mmse-smarthome}  & 22 & mmse & *** & Befehlszeile für Serververwaltung \\
    \hline
    RDP & 10.0.0.1/\textit{mmse-smarthome} & 3389 & mmse & *** & GUI-Zugriff für Wartung und Demonstration \\
    \hline
    Ubuntu Server & SSH/RDP & 22 / 3389 & mmse & *** & Benutzeraccount für Kamera-Einbundung\\
    \hline
    Home Assistant & 10.0.0.1:8123 & 8123 & mmse & *** & lokale Bedienoberfläche für Smart-Home-Logik \\
    \hline
    TP-Link & -- & -- & Johannes.Baischer@nmmse.at & *** & Zugang zur Konfiguration der Kamera\\
    \hline
    Kamera & -- & -- & mmse & *** & für Gerätedokumentation und Monitoring \\
\end{tabularx}

\section{Verwendete Dienste und lokale Adressen}

Die praktische Infrastruktur basiert auf einem klaren Muster: Das Projekt-Netz ist lokal und wird vom Server als Gateway und Router für die Lernumgebung genutzt. Dadurch bleiben die Endgeräte untereinander erreichbar, ohne dass sie aus dem Schulnetz herauslaufen müssen.

\begin{description}[leftmargin=1.8em,style=nextline]
    \item[Server-Host] Der Basisserver läuft als Xubuntu-Host mit dem Namen \texttt{mmse-smarthome}. Die lokale Verbindung nutzt dabei in der Regel die Adresse \texttt{10.0.0.1} oder den Hostnamen im selben Netzwerk.
    \item[Router] Der A1-Router bleibt als Verwaltungs- und Verteilungsgerät im lokalen Netz aktiv, aber ohne DHCP-Server (diese Aufgabe übernimmt der Server-Host). Der Router fungiert also nach seiner Konfiguration als \textbf{Access Point} Seine Adresse bleibt \texttt{10.0.0.138}.
    \item[Home Assistant] Die Automationsoberfläche läuft in der Regel auf Port \texttt{8123}. Damit ist sie auf dem lokalen Netzwerk per Browser erreichbar, zum Beispiel über \texttt{http://10.0.0.1:8123}.
    \item[SSH und RDP] Für Administration und Fernwartung werden SSH auf Port \texttt{22} sowie Remote Desktop auf Port \texttt{3389} genutzt. Beide Dienste greifen auf denselben Server-Account zu.
    \item[Tailscale] \Sidebar{Fernzugriff} Für externen Fernzugriff wird ein Overlay-VPN über Tailscale eingerichtet. Dadurch bleibt der Zugriff auch außerhalb des lokalen Netzes möglich, ohne Ports im Schulnetz freizugeben.
\end{description}

\section{Sinnvolle Arbeitsweise / Trouble Shooting}

\begin{itemize}[leftmargin=1.3em,label=\textcolor{SmartTeal}{\small\textbullet}]
    \item Kennzeichne die Zugangsdaten immer am Anfang der Projektphase, damit alle Beteiligten denselben Zugang verwenden.
    \item Nutze getrennte Benutzer für Server, Kamera und Home Assistant, um Fehlerquellen zu reduzieren.
    \item Verwende möglichst lokale Zugriffspfadnamen wie \texttt{mmse-smarthome} statt nur IP-Adressen, wenn du dich im selben Netzwerk befindest.
    \item Stelle sicher, dass alle Geräte nach dem Netzwerksetup einmal neu gestartet werden, damit sie die neue IP-Konfiguration übernehmen (besonders nach einmaligem DHCP Wechsel von Router auf Server-Host).
    \item \Sidebar{Internet Zugang} Nachdem der A1-Router als Access-Point konfiguriert wird, kann der Server sich mit einem beliebigem zweiten Netzwerk verbinden, um \textbf{Internet Zugang} zu erhalten. 
\end{itemize}
